% Final held-out numbers belong here only after the frozen evaluation completes.
\subsection{Preliminary Validation and Results Status}
Both map families completed the same two-stage training procedure. On the shared development subset, teacher-rollout refinement improved the one-step It\^o map CRPS from 0.3661 to 0.3047 and the one-step deterministic flow map from 0.3088 to 0.2885. The deterministic flow map therefore has the stronger observed one-step validation score. At eight steps, the refined map scores are closer: 0.2942 for the It\^o map and 0.2932 for the flow map. Table~\ref{tab:validation-development} reports these preliminary measurements with the teacher references and uncertainty intervals.

Both maps reduced measured sampling time. However, their quality depends on the reference budget: the one-step It\^o map improved on the 16-step SI teacher's CRPS (0.3177) while remaining worse than its 64-step score (0.2830). The eight-step It\^o score was about 3.9\% worse than SI64 with about $6.9\times$ lower measured latency on this subset. These comparisons support studying a quality--cost tradeoff; they do not establish superior stochastic-map forecasting or equality with a converged teacher.

These results use only four validation trajectories, eight ensemble members, and one training seed. They were used to guide development and must not be interpreted as held-out test performance. Final configuration selection is ongoing; held-out testing has not started. The planned final evaluation covers all 20 test trajectories with one-interval forecasts and ten-lead autoregressive forecasts, using frozen checkpoints and configurations. A controlled $K=16$ refinement is being compared with $K=4$ using the same 500-update budget; feature count will be selected by one-step validation CRPS, with one checkpoint retained for every sampling budget. Two additional complete student training seeds are planned after this selection. These future results are not yet included.

% Final layout after test evaluation:
% 1. Main table: all four model families; sampling NFE, CRPS, energy score,
%    RMSE, ensemble latency, trajectory bootstrap intervals.
% 2. CRPS-vs-NFE and CRPS-vs-latency frontier; paired model differences.
% 3. Ten physical leads: CRPS/rank/spread and spectral-error curves;
%    no teacher-forced truth inserted after initialization.
% 4. Keep this development comparison in appendix after final test narrative.
